\chapter[El Infierno de las Sombras]{El Infierno de las Sombras\\ \large El fuego que enciendes en otros, arderá en tu propia alma.}

\elegantimage{10_infierno_sombras/web/assets/hero.png}

\lettrine[lines=3, nindent=0.5em, findent=0.2em]{A}{}rrebatar la vida o el sustento de otro es una deuda que no se salda fácilmente. Es introducir en el universo una fuerza de destrucción que, por ley natural, debe regresar con la misma intensidad.

\section{El Vestido de Fuego}
\begin{elegantquote}
\textit{Había un hombre implacable que construyó su poder sobre la sangre y el llanto de los inocentes. Creía que la crueldad era un escudo que nadie podría traspasar. Se reía de las lágrimas y no sentía ninguna piedad en el corazón...}

\textit{Pero la vida es como un gran eco. Todo lo que lanzamos al abismo vuelve a nosotros con la misma fuerza vibratoria. Por eso, al cerrar los ojos, no encontró la nada; se halló atrapado en un valle de un calor asfixiante y un terror profundo. No eran castigos inventados ni demonios ajenos...}

\textit{Eran todas y cada una de las agonías, miedos y gritos de quienes asesinó, concentrados alrededor suyo como un fuego abrasador que le gritaba la verdad. A través de este calor insoportable, la roca de su insensibilidad se fue fundiendo... hasta que, arrodillado frente a su propio dolor, suplicó el don de la compasión.}

\end{elegantquote}

\vspace{1em}
\begin{center}
    \textbf{\Large Nada se destruye, solo vuelve a su origen.}
\end{center}
\vspace{1em}

\section{Análisis Kármico y Traducción}
\begin{wrapfigure}{r}{0.5\textwidth}
  \vspace{-1.5em}
  \begin{center}
  \inlineelegantimage[0.45\textwidth]{10_infierno_sombras/web/assets/pasaje_original.jpg}
  \end{center}
  \vspace{-1.5em}
\end{wrapfigure}
Quien siembra dolor extremo, robo y asesinato, atrae hacia sí mismo un sufrimiento ardiente y cortante, un verdadero infierno personal diseñado para disolver su insensibilidad de piedra.

